\section{Parte III. Vecinos cercanos}
\label{sec:parte_03}

\subsection*{Inciso 8. KNN con distintos valores de K}

Se procede a estimar vecinos cercanos considerando los siguientes valores del hiperparámetro $K={10,50,100}$. Una vez estimados los modelos para los distintos valores correspondientes de $K$, se calcula el porcentaje de observaciones que el modelo clasifica correctamente \textit{()accuracy)}. Los resultados que se obtienen son los siguientes:

\input{\joinpath{\outputspath}{B_Table/tabla_knn_k_fijos.tex}}


La elección de $K$ introduce un \textit{trade-off} entre sesgo y varianza, dado que a medida que aumenta el número de vecinos cercanos se consideran observaciones cada vez más alejadas de la observación a clasificar y por lo tanto menos similares, lo que en general reduce la varianza pero aumenta el sesgo. En el extremo cuando $K=1$, el modelo considera únicamente al vecino más cercano, por lo que presenta un sesgo bajo y una varianza elevada, ya que pequeños cambios en los datos pueden modificar la clasificación. En cambio, cuando $K=n$, se utiliza toda la muestra de entrenamiento para clasificar según la mayoría, lo que reduce la varianza pero aumenta el sesgo dado que considero a todas estas observaciones y no necesariamente a las que están más cercanas.


\subsection*{Inciso 9. Selección de K cross validation}


Se selecciona el valor óptimo de $K$ mediante cross validation utilizando una division de cinco partes de la muestra de entrenamiento. De esta manera, en cada iteración se utilizan cuatro partes para estimar el modelo y la restante para evaluar su capacidad de clasificación, calculando el \textit{accuracy}. Luego este procedimiento se repite para distintos valores de $K$, entre 1 y 1000, y se obtiene el \textit{accuracy} promedio de los cinco bloques para cada valor.

Para evitar utilizar información del bloque de validación, los valores faltantes se imputan y las variables se estandarizan utilizando únicamente las observaciones correspondientes al entrenamiento de cada bloque.

\input{\joinpath{\outputspath}{B_Table/tabla_knn_cv.tex}}

\begin{figure}[H]
	\centering
	\includegraphics[width=0.95\textwidth]{\joinpath{\outputspath}{B_Figure/knn_accuracy_cv.png}}
	\caption{Selección del número de vecinos mediante validación cruzada}
	\label{fig:knn_cv}
\end{figure}

Como se observa en la tabla y el gráfico, el \textit{accuracy} promedio alcanza su mayor valor cuando $K=300$, con un 77,40\% de observaciones correctamente clasificadas. Para valores mayores de $K$, el \textit{accuracy} comienza a disminuir, por lo que se selecciona $K=300$ como el valor óptimo dentro de la grilla evaluada (KNN con K-CV).
